% !TEX root = ../main.tex
\chapter{Introduction}\label{ch:introduction}

% Funnel from the broad context to the specific problem. Typical parts:
% context, problem statement and research gap, aim and research questions,
% contributions, scope, and an outline of the thesis.

\section{Context and motivation}\label{sec:context}

% TODO: write the context. Every claim taken from the literature needs a citation.
Typesetting systems such as \LaTeX{} separate content from presentation
\autocite{lamport1994latex}, an idea related to literate programming
\autocite{knuth1984literate}. Documents are exchanged as \gls{pdf} files and
cited works are located through their \gls{doi}.

\section{Problem statement and research questions}\label{sec:problem}

\section{Contributions}\label{sec:contributions}

\section{Thesis outline}\label{sec:outline}

\Cref{ch:background} reviews the literature, \cref{ch:methodology} describes the
method, \cref{ch:results} reports the results, \cref{ch:discussion} discusses
them and \cref{ch:conclusion} concludes.
